\begin{textsample}{NEG Dim 3 – vem\_ed\_20 – Score 3.00 – vem\_ed\_20/t102.txt}  \label{ex:f3_neg_117}
continuação Desde 2021, Silvio Ribeiro desenvolve PCIs na Fatec Lins com a Universidade de Aveiro (Portugal).
Sua experiência mais recente envolve as comunidades das cidades de Lins (SP) e Murtosa (Portugal).
Professores e alunos da Fatec Lins e da Universidade de Aveiro preparam um Guia Turístico Eletrônico para as duas cidades.
No segundo semestre de 2022, em contato com a Secretaria Municipal de Cultura de Lins, foi possível abrir duas vagas de estágio na Prefeitura para estudantes da Fatec Lins contribuírem com o levantamento de dados para produção desse guia, que está em fase de protótipo.
Futuramente, o modelo poderá ser replicado pela cidade de Murtosa. “Ensino de língua estrangeira e Intercâmbios Virtuais” pautou o diálogo entre Vitor Ierusalimschy (chefe do setor de Projetos Educacionais na Superintendência de Relações Internacionais na Universidade Federal Fluminense) e Claudia Delfino (professora da Fatec Praia Grande).
Vitor Ierusalimschy destacou a importância de incluir elementos de consciência linguística e imperialismo linguístico no planejamento de projetos COIL (Collaborative Online International Learning). “É necessário abordar a situação de vulnerabilidade do falante não nativo”, ressaltou.
Para isso, é importante criar um ambiente confortável e seguro, conversar previamente com os alunos, criar atividades de interação, promover a tolerância por parte dos falantes nativos e introduzir a língua nativa do parceiro em falas iniciais.
A melhoria da proficiência em inglês na interação com falantes não nativos é o mote do PCI desenvolvido por Claudia Delfino na Fatec Praia Grande.
Neste projeto encontram-se: Sapir Academic College e Kaye Academic College, de Israel; University of Opole, da Polônia e Universitá Cattolica del Sacro Cuore, da Itália.
A professora destaca que, com a participação dos alunos em PCIs, nota-se uma melhora perceptível na produção linguística dos participantes."A língua é só uma isca.
O principal é a convivência com a diversidade cultural, incluindo beduínos israelenses, o desenvolvimento das habilidades de resolução de conflitos e de trabalho em equipes internacionais".

% matched lemmas: 
\end{textsample}
